% titel: Fluchtgeschwindigkeit
% skript: Kap. Gravitation, Energien im Gravitationsfeld

Eine auf der Erdoberfläche abgeschossene Kugel soll das Gravitationspotential
der Erde verlassen. Massgebend ist die Energiebilanz im Newtonschen
Zentralfeld: am Anfang kinetische plus (negative) potentielle Energie, im
Unendlichen ist die potentielle Energie null.

\begin{werkzeugkasten}[Werkzeuge]
\textbf{Gravitationskraft} (Punktmassen/radialsymmetrischer Körper):
\begin{equation}
  F_G = G\,\frac{M_E\, m}{r^2}\;.
\end{equation}

\textbf{Potentielle Energie} (Nullpunkt im Unendlichen):
\begin{equation}
  E_\mathrm{pot} = -G\,\frac{M_E\, m}{r}\;.
\end{equation}

An der Oberfläche $r=R_E$ ist $E_\mathrm{pot}=-G\frac{M_E m}{R_E}$.

\textbf{Energieerhaltung} bei reibungsfreier, vertikaler Bewegung ohne äussere
Kräfte:
\begin{equation}
  \frac12 m v^2 - G\frac{M_E m}{r} = \mathrm{konst}.
\end{equation}

\textbf{Fluchtbedingung:} Um ins Unendliche ($r\to\infty$) zu gelangen, muss
für $r\to\infty$ noch $v\ge 0$ (kinetische Energie nicht negativ) sein.
Diese Bedingung ergibt die Fluchtgeschwindigkeit.

\textbf{Konstante:} $G=\qty{6.674e-11}{\metre\cubed\per\kilogram
\per\square\second}$; $M_E=\qty{5.974e24}{\kg}$, $R_E=\qty{6378}{km}$.
\end{werkzeugkasten}

\begin{werkzeugkasten}[Vorgehen in allgemeiner Form]
Energiebilanz zwischen Start ($r=R_E$, $v=v_0$) und Unendlich
($r\to\infty$, $v\to 0$)
aufstellen und nach $v_0$ auflösen.
\end{werkzeugkasten}

\begin{achtung}[Typische Fallen]
\begin{itemize}
  \item Die Masse der Kugel kürzt sich in der Energiebilanz vollständig
       heraus: die Fluchtgeschwindigkeit ist massenunabhängig.
  \item Das Vorzeichen der potentiellen Energie (hier negativ) unbedingt
       beibehalten.
\end{itemize}
\end{achtung}
